\chapter{Recursion}

\emph{Recursion} is a term used to refer to something defined in terms
of itself. An example is the factorial: \(n! = (n - 1)!\), with \(0! = 1\) being the
base case where the recursion stops.

\section{From Recursive to Closed Form}

Often, a sequence defined recursively can be rewritten in an explicit or closed form.

\emph{Recursive form} defines \( a_n \) based on previous terms (e.g., \( a_n = a_{n-1} + 2 \)).

\emph{Closed form} expresses \( a_n \) directly in terms of \(n\), with no reference to previous terms.

\section{Linear Recursion}

Linear recursion has the form:

\begin{align*}
    a_n & = \lambda a_{n -1} + c \\
    a_1 & = \lambda a_0 + c \implies a_2 = \lambda (\lambda a_0 + c) + c \\
    \implies a_n & = \lambda^n a_0 + (1 + \lambda + \lambda^2 + \cdots + \lambda^{n - 1})c.
\end{align*}

Therefore,

\[
    a_n = \lambda^n a_0 + \frac{\lambda^n - 1}{\lambda - 1}c.
\]

\subsection{General Method for Linear Recursion}

\begin{enumerate}

    \item Calculate some terms and observe the difference or quotient between them.
    Maybe the pattern is clear, and you can build a geometric or arithmetic sequence directly.

    \item Identify the type (e.g., linear, homogeneous, constant coefficients).

    \item Solve the characteristic equation (if linear).

    \item Use initial conditions to determine constants.

\end{enumerate}


\section{Recursive Inequalities}

\emph{Recursive inequalities} bound a sequence rather than define it exactly.

\textbf{Example:}

Suppose

\[
    a_{n+1} \le \frac{1}{2} a_n + 3.
\]

To analyze

\begin{itemize}

    \item Guess a bound (e.g., show \( a_n \le M \) for all \(n\)).
    \item Use induction or iteration to justify convergence or boundedness.
    \item Compare with a simpler sequence (e.g., geometric series).

\end{itemize}

These techniques are common in analysis, approximation, and numerical algorithms.

\section{Closed Form of the Fibonacci Sequence}

The Fibonacci sequence is defined recursively:

\[
    F_0 = 0, \quad F_1 = 1, \quad F_n = F_{n-1} + F_{n-2}.
\]

The characteristic equation is:

\[
    x^2 - x - 1 = 0 \Rightarrow x = \varphi = \frac{1 + \sqrt{5}}{2}, \quad \psi = \frac{1 - \sqrt{5}}{2}.
\]

The closed form (Binet’s Formula) is:

\[
    F_n = \frac{1}{\sqrt{5}} \left( \varphi^n - \psi^n \right).
\]

\begin{itemize}
    
    \item \( \varphi \approx 1.618 \) is the golden ratio.
    
    \item This formula demonstrates exponential growth of \( F_n \).

\end{itemize}

This formula can be derived using the steps of the chapter about \textbf{generating functions}.

\section{Recurrent Sequences and Strong Induction}

To prove the \emph{closed form} of a given recurrence, we use the principle of strong induction and the inductive definition.

\textbf{Example:}

Given is the recurrence relation 

\[
    c_k = 3c_{k - 1} + 1, \quad c_1 = 1
\]

with the explicit formula 

\[
    c_n = \frac{3^n - 1}{2}, \quad \forall n \ge 1
\]

\textbf{Proof:}

Let \(n = 1\), then \(c_1 = 1\).

We will assume our closed formula for some \(k\), \(1 \le j \le k\) and we want to prove that 

\[
    c_{n + 1} = \frac{3^{n + 1} - 1}{2}
\]

Since our formula applies for all natural numbers between 1 and \(k\), we can use our recurrence relation and show that 

\begin{align*}
    c_k = 3c_{k - 1} + 1 &= 3 \left(\frac{3^n - 1}{2}\right) + 1 \\ 
    &= \frac{3^{n + 1} - 3}{2} + 1 \\ 
    &= \frac{3^{n + 1} - 3 + 2}{2} \\ 
    &= \frac{3^{n + 1} - 1}{2}
\end{align*}

Hence, we have proven our closed form.

\QED